\section{Conclusion}
The analysis demonstrates that the experimental gamma-ray spectrum contains seven prominent spectral features that can be identified and characterised through peak detection and Gaussian fitting. The fitted peak centres provide a consistent set of experimental positions, while comparison with the reference gamma-ray energies allows the detector's energy calibration to be established. The resulting calibration is well represented by a linear relationship, with a slope of approximately 1.19174 keV per channel and an RMSE of approximately 2.90 keV.

The agreement between the calibrated energies and reference transition energies indicates that the selected peak assignments are internally consistent. The analysis also shows that the widths of the spectral features vary across the measured energy range, with some peaks exhibiting broader structures that may involve unresolved contributions.

Overall, the obtained calibration provides a basis for converting the measured spectrum from its uncalibrated coordinate into physical gamma-ray energy. The calibrated spectrum can subsequently be used for identifying gamma-ray transitions and performing further quantitative analysis of the detector response.